\documentclass[11pt,a4paper]{article}

\usepackage[utf8]{inputenc}
\usepackage[T1]{fontenc}
\usepackage{lmodern}
\usepackage[margin=1in]{geometry}
\usepackage{amsmath,amssymb,amsthm,mathtools}
\usepackage{booktabs}
\usepackage{array}
\usepackage{enumitem}
\usepackage{xcolor}
\usepackage{listings}
\usepackage{hyperref}

\hypersetup{
  colorlinks=true,
  linkcolor=blue!60!black,
  urlcolor=blue!60!black,
  citecolor=blue!60!black
}

\lstdefinelanguage{Lean4}{
  keywords={def, theorem, lemma, by, have, let, in, if, then, else, match, with, fun, intro, rcases, rw, exact, refine, apply, omega, decide, simp, dsimp, ext, subst, generalize, constructor, forall, exists, import, namespace, end, open, use},
  keywordstyle=\color{blue!70!black}\bfseries,
  comment=[l]{--},
  morecomment=[s]{/-}{-/},
  commentstyle=\color{green!45!black}\itshape,
  basicstyle=\ttfamily\small,
  breaklines=true,
  columns=fullflexible,
  keepspaces=true,
  frame=single,
  rulecolor=\color{black!20},
  backgroundcolor=\color{black!2},
  literate=
    {ℕ}{{$\mathbb{N}$}}1
    {ℤ}{{$\mathbb{Z}$}}1
    {×}{{$\times$}}1
    {∀}{{$\forall$}}1
    {∃}{{$\exists$}}1
    {∈}{{$\in$}}1
    {≤}{{$\le$}}1
    {≥}{{$\ge$}}1
    {∧}{{$\land$}}1
    {∨}{{$\lor$}}1
    {→}{{$\to$}}1
    {⇒}{{$\Rightarrow$}}1
    {⟨}{{$\langle$}}1
    {⟩}{{$\rangle$}}1
}

\newtheorem{theorem}{Theorem}[section]
\newtheorem{lemma}[theorem]{Lemma}
\newtheorem{proposition}[theorem]{Proposition}
\theoremstyle{remark}
\newtheorem{remark}[theorem]{Remark}

\title{\textbf{Formalizing the Generalized $n$-Queens Problem in Lean 4:\\
An Audit and Completion of the Le--Li--Wang (1990) Proof}}
\author{Luc Duponcheel \\ \small Formalized in Lean 4 + Mathlib (\texttt{n\_c\_coins\_placement\_proof.lean})}
\date{October 2026}

\begin{document}

\maketitle

\begin{abstract}
We report on the complete, \texttt{sorry}-free and \texttt{native\_decide}-free formalization in Lean~4 (using Mathlib) of the Generalized $n$-Queens Conjecture (MathOverflow~510135): for every board size $n \ge 4$ and every integer $0 \le c \le n$, one can place $nc$ coins on an $n \times n$ chessboard (at most one per cell) such that every row, column, downward diagonal, and upward diagonal contains at most $c$ coins. Our formalization follows the constructive method of Le, Li, and Wang (1990), \emph{``A note on the generalized $n$-queens problem''}. Along the way, formal verification exposed several significant gaps, unstated boundary phenomena, and delicate modular cancellations in the 2-page paper that had to be supplied or corrected to obtain a complete proof.
\end{abstract}

\tableofcontents

\section{Introduction and Problem Statement}

In MathOverflow question~510135 (\href{https://mathoverflow.net/questions/510135}{\texttt{mathoverflow/510135}}), the following generalization of the classical $n$-queens problem was posed:

\begin{theorem}[Generalized $n$-Queens / $(n,c)$-Coins Placement]
\label{thm:main}
For every integer $n \ge 4$ and every integer $0 \le c \le n$, there exists a subset of cells
\[
  S \subseteq \{0, 1, \dots, n-1\}^2
\]
of cardinality $|S| = nc$ such that:
\begin{enumerate}[nosep,label=(\roman*)]
  \item \textbf{Row bound:} For every $0 \le i < n$, $|\{(r, \mathrm{col}) \in S \mid r = i\}| \le c$.
  \item \textbf{Column bound:} For every $0 \le j < n$, $|\{(r, \mathrm{col}) \in S \mid \mathrm{col} = j\}| \le c$.
  \item \textbf{Downward-diagonal bound:} For every $0 \le d < 2n - 1$, $|\{(r, \mathrm{col}) \in S \mid r + \mathrm{col} = d\}| \le c$.
  \item \textbf{Upward-diagonal bound:} For every $0 \le d < 2n - 1$, $|\{(r, \mathrm{col}) \in S \mid (n - 1) + r - \mathrm{col} = d\}| \le c$.
\end{enumerate}
\end{theorem}

When $c = 1$, Theorem~\ref{thm:main} reduces to the classical $n$-queens problem ($n \ge 4$). Because $|S| = nc$ and the $n$ rows (resp.\ $n$ columns) partition the board, the row and column inequalities $\le c$ are in fact equalities ($= c$ coins per row and per column), whereas the $2n - 1$ diagonals in each direction have lengths ranging from $1$ to $n$ and therefore satisfy an inequality $\le c$.

In a 2-page note published in 1990~\cite{LLW1990}, Mao-Hua Le, Kang-Di Li, and Zhen-Fa Wang gave a constructive proof of Theorem~\ref{thm:main}. We have formalized the complete proof of Theorem~\ref{thm:main} in Lean~4 (\texttt{n\_c\_coins\_placement\_proof.lean}, $\approx 4{,}450$ lines, importing only \texttt{Mathlib}, with zero \texttt{sorry}s and zero \texttt{native\_decide}s).

This note summarizes:
\begin{enumerate}[nosep]
  \item The structure of the Lean~4 formalization;
  \item What parts of the Le--Li--Wang (LLW) construction work verbatim; and
  \item What mathematical gaps and edge cases in~\cite{LLW1990} were uncovered during formalization and how they were resolved.
\end{enumerate}

\section{Overview of the Lean 4 Formalization}

In \texttt{n\_c\_coins\_placement\_proof.lean}, a placement on the $n \times n$ board is represented as a finite set \texttt{S : Finset ($\mathbb{N} \times \mathbb{N}$)}, and the validity predicate \texttt{n\_c\_placement\_conjecture\_prop n c S} bundles the six conditions:
\begin{lstlisting}[language=Lean4]
def f_cells_ok_prop (n : ℕ) (S : Finset (ℕ × ℕ)) : Prop :=
  ∀ p ∈ S, p.1 < n ∧ p.2 < n

def f_rows_ok_prop (n c : ℕ) (S : Finset (ℕ × ℕ)) : Prop :=
  ∀ i, i < n → #(S.filter (fun p => p.1 = i)) ≤ c

def f_cols_ok_prop (n c : ℕ) (S : Finset (ℕ × ℕ)) : Prop :=
  ∀ j, j < n → #(S.filter (fun p => p.2 = j)) ≤ c

def f_dds_ok_prop (n c : ℕ) (S : Finset (ℕ × ℕ)) : Prop :=
  ∀ d, d < 2 * n - 1 → #(S.filter (fun p => p.1 + p.2 = d)) ≤ c

def f_uds_ok_prop (n c : ℕ) (S : Finset (ℕ × ℕ)) : Prop :=
  ∀ d, d < 2 * n - 1 → #(S.filter (fun p => (n - 1) + p.1 - p.2 = d)) ≤ c

def f_bc_ok_prop (n c : ℕ) (S : Finset (ℕ × ℕ)) : Prop :=
  #S = n * c
\end{lstlisting}

The master construction \texttt{llw\_board n c} dispatches on $c$ into four regimes:
\begin{enumerate}[nosep]
  \item $c = 0$: the empty board $\emptyset$.
  \item $c = 1$: an explicit classical $n$-queens configuration \texttt{classic\_n\_queens n}.
  \item $c \in \{n - 2, n - 1, n\}$: a complement construction \texttt{simple\_case n c} obtained by removing $0$, $1$, or $2$ disjoint transversals from the full $n \times n$ board.
  \item $2 \le c \le n - 3$: the Le--Li--Wang modular decomposition
  \[
    c = k_0 + 6t, \qquad k_0 = 2 + ((c - 2) \bmod 6) \in \{2, 3, 4, 5, 6, 7\}, \qquad t = \frac{c - k_0}{6},
  \]
  which takes the disjoint union of a base board \texttt{construct\_base n k0} of size $nk_0$ and $t$ cyclic strips \texttt{llw\_F n 6 (k0 + 2 + 6 * e)} for $0 \le e < t$.
\end{enumerate}

The final theorem in \texttt{n\_c\_coins\_placement\_proof.lean} is completely proved:
\begin{lstlisting}[language=Lean4]
theorem n_c_coins_placement_conjecture :
  ∀ (n : ℕ) (c : ℕ) (_hn : n ≥ 4) (_hc : c ≤ n),
    ∃ board : Finset (ℕ × ℕ),
    n_c_placement_conjecture_prop n c board := by
  intro n c hn hc
  use llw_board n c
  exact ⟨
    llw_board_cells_ok n c hn hc,
    llw_board_rows_ok n c hn hc,
    llw_board_cols_ok n c hn hc,
    llw_board_dds_ok n c hn hc,
    llw_board_uds_ok n c hn hc,
    llw_board_bc_ok n c hn hc
  ⟩
\end{lstlisting}

\section{The Core Le--Li--Wang Architecture}

The key building block of~\cite{LLW1990} (translated to $0$-based indexing) is the cyclic slope-$2$ strip of width $h$ starting at column offset $s$:
\[
  F(n, h, s) = \bigl\{\, \bigl(r,\, (2r + s + t) \bmod n\bigr) \;\big|\; 0 \le r < n,\; 0 \le t < h \,\bigr\},
\]
formalized as \texttt{llw\_F n h s}.

Why does $F(n, h, s)$ almost solve the problem by itself, and where does it fail?
\begin{itemize}
  \item \textbf{Rows (\texttt{llw\_F\_rows\_ok}) and Cardinality (\texttt{llw\_F\_bc\_ok}):} For $h \le n$, each row $r$ contains $h$ distinct columns $(2r + s + t) \bmod n$, so $|F(n, h, s)| = nh$ and each row has exactly $h$ coins.
  \item \textbf{Upward diagonals (\texttt{llw\_F\_uds\_ok}):} For a fixed shift $t$, if two cells $(r_1, c_1)$ and $(r_2, c_2)$ lie on the same upward diagonal, then $r_1 - c_1 = r_2 - c_2$ in $\mathbb{Z}$. Since $c_i \equiv 2r_i + s + t \pmod n$, we obtain
  \[
    r_1 - (2r_1 + s + t) \equiv r_2 - (2r_2 + s + t) \pmod n \implies r_1 \equiv r_2 \pmod n,
  \]
  which forces $r_1 = r_2$ for all $n$ (with no coprimality assumption on $n$!). Thus each of the $h$ shifts contributes at most $1$ coin to any upward diagonal, giving $\le h$ coins in total.
  \item \textbf{Columns (\texttt{llw\_F\_cols\_ok}):} For a fixed shift $t$, $(2r + s + t) \bmod n = j$ has:
  \begin{itemize}[nosep]
    \item $1$ solution $r$ when $n$ is odd ($\gcd(2, n) = 1$);
    \item $2$ solutions $r$ and $r + n/2$ when $n$ is even and $j \equiv s + t \pmod 2$, and $0$ solutions when $j \not\equiv s + t \pmod 2$.
  \end{itemize}
  Consequently, $F(n, h, s)$ satisfies the column bound $\le h$ whenever \textbf{$n$ is odd or $h$ is even} (since consecutive shifts $t$ and $t+1$ have opposite parities and share the $2$ solutions per pair of shifts). However, \textbf{when $n$ is even and $h$ is odd, $F(n, h, s)$ violates the column bound} (having $h + 1$ coins in columns matching the parity of $s$).
  \item \textbf{Downward diagonals (\texttt{llw\_F\_three\_dds\_ok}, \texttt{llw\_F\_dds\_not\_div3}):} A cell $(r, c) \in F(n, h, s)$ with $r + c = d$ satisfies one of the three wrap-around equations:
  \begin{equation}
  \label{eq:wrap}
    3r + s + t = d, \qquad 3r + s + t = d + n, \qquad 3r + s + t = d + 2n.
  \end{equation}
  \begin{itemize}[nosep]
    \item When $3 \nmid n$, for a fixed shift $t$, the right-hand sides $d, d + n, d + 2n$ cover all three residue classes modulo $3$, so at most \emph{one} of the three equations in~\eqref{eq:wrap} can be divisible by $3$. Hence each shift $t$ contributes at most $1$ coin to diagonal $d$, and $F(n, h, s)$ has $\le h$ coins per downward diagonal for \emph{any} $h$ (\texttt{llw\_F\_dds\_not\_div3}).
    \item When $3 \mid n$, however, $d \equiv d + n \equiv d + 2n \pmod 3$. Thus if $s + t \equiv d \pmod 3$, a \emph{single} shift $t$ can have up to \textbf{3 coins} on the same downward diagonal $d$ (one in each of the three wrap-around intervals)! Conversely, if $h$ is a multiple of $3$ ($h = 3$ or $h = 6$), only one out of every three consecutive shifts satisfies $s + t \equiv d \pmod 3$, so $F(n, 3, s)$ still has $\le 3$ coins per downward diagonal (\texttt{llw\_F\_three\_dds\_ok}) and $F(n, 6, s)$ has $\le 6$ coins per downward diagonal (\texttt{llw\_F\_six\_dds\_ok}).
  \end{itemize}
\end{itemize}

This explains why LLW decompose $c = k_0 + 6t$: a block $F(n, 6, s)$ of width $6$ has \emph{even} width (so columns have $\le 6$ coins for all $n$) and width \emph{divisible by $3$} (so downward diagonals have $\le 6$ coins even when $3 \mid n$). All of the difficulty is concentrated in the base cases $k_0 \in \{2, 3, 4, 5, 6, 7\}$ when $2 \mid n$ (for odd $k_0 \in \{3, 5, 7\}$) or $3 \mid n$ (for $k_0 \in \{4, 7\}$).

\section{Audit of the Le--Li--Wang (1990) Proof: Gaps and Corrections}

While the $c = k_0 + 6t$ decomposition and the perturbed strips $T_1, T_2, T_4, T_5$ in~\cite{LLW1990} are genuinely insightful, formalizing the paper in Lean~4 revealed three substantive issues that a careful referee report must highlight.

\subsection{Gap 1: Modular Wrap-Around Overflow at $c = n - 2$ and $c = n - 1$}

In~\cite{LLW1990}, the perturbed strips $T_1(n, s)$ and $T_2(n, s)$ occupy shifts $s-1$ and $s$, and $T_4(n), T_5(n)$ occupy shifts up to $k_0 + 1$. Because of this, the $t$ extension blocks of width $6$ cannot start at shift $k_0$; they must start at shift $s = k_0 + 2$:
\[
  \bigcup_{e = 0}^{t - 1} F(n, 6, k_0 + 2 + 6e).
\]
Consequently, the shifts used by the construction span from $0$ up to $k_0 + 2 + 6t - 1 = c + 1$. For these shifts to remain disjoint from the initial columns of $F(n, k_0, 0)$ modulo $n$, one requires:
\begin{itemize}[nosep]
  \item $c + 2 \le n$ (i.e., $c \le n - 2$) whenever $t \ge 1$, and
  \item $k_0 + 1 \le n$ (or $k_0 + 2 \le n$ for $T_4, T_5$) whenever $t = 0$.
\end{itemize}
When $c = n - 1$ or $c = n - 2$, the LLW formula wraps around modulo $n$ and places multiple coins on the same cell, violating $|S| = nc$.

\paragraph{Resolution in Lean 4 (\texttt{simple\_case}):}
For $c \in \{n - 2, n - 1, n\}$, we use a direct complement construction from the full board $B = \{0, \dots, n-1\}^2$:
\begin{itemize}
  \item \textbf{$c = n$:} $S = B$.
  \item \textbf{$c = n - 1$:} We remove a single $n$-transversal consisting of the two anti-diagonal corners $(0, n-1), (n-1, 0)$ and the interior of the main diagonal $\{(i, i) \mid 1 \le i \le n - 2\}$. Every row and column loses $1$ coin, and the two longest diagonals ($r + \mathrm{col} = n - 1$ and $r = \mathrm{col}$, which are the only diagonals of length $n > n - 1$) lose $2$ and $n - 2 \ge 2$ coins respectively.
  \item \textbf{$c = n - 2$ ($n$ even):} We remove both the main diagonal $D_1 = \{(i, i)\}$ and the anti-diagonal $D_2 = \{(i, n - 1 - i)\}$. Since $n$ is even, $D_1 \cap D_2 = \emptyset$, so $|B \setminus (D_1 \cup D_2)| = n(n - 2)$, every row and column loses $2$ coins, and all diagonals of length $n$ or $n - 1$ have $\le n - 2$ coins remaining.
  \item \textbf{$c = n - 2$ ($n$ odd):} When $n = 2m + 1$ is odd, $D_1$ and $D_2$ intersect at the center cell $(m, m)$, so $|D_1 \cup D_2| = 2n - 1 \ne 2n$. In \texttt{simple\_case}, we locally reroute the two removed transversals around the center $m = (n - 1)/2$:
  \begin{align*}
    P_1 &= \{(i, i) \mid i \notin \{m, m+1\}\} \cup \{(m, m+1), (m+1, m)\}, \\
    P_2 &= \{(i, n - 1 - i) \mid i \notin \{m-1, m, m+1\}\} \cup \{(m-1, m), (m, m-1), (m+1, m+1)\}.
  \end{align*}
  We prove in \texttt{simple\_case\_bc\_odd\_rem\_card} that $P_1 \cap P_2 = \emptyset$ and $|P_1 \cup P_2| = 2n$, while every row, column, and diagonal of length $\ge n - 1$ intersects $P_1 \cup P_2$ in at least $2$ cells.
\end{itemize}

\subsection{Gap 2: Downward-Diagonal Collisions and Cancellations when $3 \mid n$}

In~\cite{LLW1990}, the perturbed sets $T_1(n, s)$, $T_2(n, s)$ (for even $n$ and odd $s \in \{3, 5, 7\}$) and $T_4(n), T_5(n)$ (for $n \equiv 3 \pmod 6$) are introduced to fix the column/diagonal bounds of $F(n, k_0, 0)$, and the verification of their diagonal bounds is dismissed as ``easy to verify.''

During formalization, we found that \textbf{$T_1(n, s)$, $T_2(n, s)$, and $T_5(n)$ do not satisfy a $1$-coin-per-downward-diagonal bound when $3 \mid n$}; each of them can place \textbf{2 coins} on the same downward diagonal! Conversely, when $k_0 = 5$ or $k_0 = 7$, adding $2$ coins to a strip of width $4$ or $6$ would exceed $k_0$ unless the base strip has \emph{strictly fewer} than $k_0 - 1$ coins on those diagonals.

The formal proof reveals a remarkable, unstated cancellation centered on the main anti-diagonal $d = n - 1$ (and its neighbors $d = n - 2, n + 1$):

\begin{lemma}[Sharp Downward-Diagonal Bounds for $F(n, h, 0)$, Formalized in Lean 4]
\label{lem:sharp-F}
Let $n \ge 6$ and $0 \le d < 2n - 1$.
\begin{enumerate}[nosep,label=(\alph*)]
  \item \texttt{llw\_F\_two\_zero\_dds\_ok}: $|F(n, 2, 0) \cap \{r + c = d\}| \le 2$ for all $d$ (even when $3 \mid n$, because $2r + t < 2n$ for $t \in \{0, 1\}$, preventing the third wrap-around).
  \item \texttt{llw\_F\_two\_zero\_dds\_ok\_eq\_zero}: If $3 \mid n$, then on the main anti-diagonal $d = n - 1$,
  \[
    |F(n, 2, 0) \cap \{r + c = n - 1\}| = 0,
  \]
  because $3r + t \equiv t \in \{0, 1\} \pmod 3$ whereas $n - 1 \equiv 2n - 1 \equiv 2 \pmod 3$.
  \item \texttt{llw\_F\_three\_zero\_dds\_ok\_ne}: For \emph{every} $d \ne n - 1$,
  \[
    |F(n, 3, 0) \cap \{r + c = d\}| \le 2,
  \]
  because the third wrap-around $2r + t \ge 2n$ in $F(n, 3, 0)$ occurs uniquely at $(r, t) = (n - 1, 2)$, which is the cell $(n - 1, 0)$ on diagonal $d = n - 1$.
  \item \texttt{llw\_F\_six\_zero\_dds\_ok\_ne}: Since $F(n, 6, 0) = F(n, 3, 0) \cup F(n, 3, 3)$, for \emph{every} $d \ne n - 1$,
  \[
    |F(n, 6, 0) \cap \{r + c = d\}| \le 2 + 3 = 5,
  \]
  whereas for $d = n - 1$, $|F(n, 6, 0) \cap \{r + c = n - 1\}| \le 6$.
\end{enumerate}
\end{lemma}

Lemma~\ref{lem:sharp-F} is the exact mechanism that makes \texttt{construct\_base\_dds\_ok} work across all subcases of $k_0 \in \{3, 4, 5, 7\}$, as summarized in Table~\ref{tab:cancellation}.

\begin{table}[ht!]
\centering
\footnotesize
\renewcommand{\arraystretch}{1.25}
\begin{tabular}{@{}llccc@{}}
\toprule
\textbf{Case ($k_0, n \bmod 6$)} & \textbf{Base Construction} & \textbf{Diag.\ $d = n - 1$} & \textbf{Diag.\ $d = n - 2, n + 1$} & \textbf{Other $d$} \\
\midrule
$k_0 = 3,\; n \equiv 0 \pmod 6$ & $F(n,2,0) \cup T_2(n,3)$ & $0 + 2 \le 3$ & $2 + 0 \le 3$ & $2 + 1 = 3$ \\
$k_0 = 4,\; n \equiv 0 \pmod 6$ & $F(n,2,0) \cup T_2(n,3) \cup T_2(n,5)$ & $0 + 2 + 1 \le 4$ & $2 + 0 + 2 = 4$ & $2 + 1 + 1 = 4$ \\
$k_0 = 4,\; n \equiv 3 \pmod 6$ & $F(n,3,0) \cup T_4(n)$ & $3 + 1 = 4$ & $2 + 1 \le 4$ & $2 + 1 \le 4$ \\
$k_0 = 5,\; n \equiv 0 \pmod 6$ & $F(n,4,0) \cup T_1(n,5)$ & $(0 + 3) + 2 = 5$ & $\subseteq F(n,6,0)\; (\le 5)$ & $\subseteq F(n,6,0)\; (\le 5)$ \\
$k_0 = 5,\; n \equiv 2 \pmod 6$ & $F(n,4,0) \cup T_1(n,5)$ & $4 + 1 = 5$ & $\subseteq F(n,6,0)\; (\le 5)$ & $\subseteq F(n,6,0)\; (\le 5)$ \\
$k_0 = 5,\; n \equiv 4 \pmod 6$ & $F(n,4,0) \cup T_2(n,5)$ & $4 + 1 = 5$ & $\subseteq F(n,6,0)\; (\le 5)$ & $\subseteq F(n,6,0)\; (\le 5)$ \\
$k_0 = 7,\; n \equiv 0,4 \pmod 6$ & $F(n,6,0) \cup T_1(n,7)$ & $6 + 1 = 7$ & $5 + 2 = 7$ & $5 + 2 = 7$ \\
$k_0 = 7,\; n \equiv 2 \pmod 6$ & $F(n,6,0) \cup T_2(n,7)$ & $6 + 1 = 7$ & $5 + 2 = 7$ & $5 + 2 = 7$ \\
$k_0 = 7,\; n \equiv 3 \pmod 6$ & $F(n,6,0) \cup T_5(n)$ & $6 + 1 = 7$ & $5 + 2 = 7$ & $5 + 2 = 7$ \\
\bottomrule
\end{tabular}
\caption{Exact downward-diagonal budgets in \texttt{construct\_base\_dds\_ok} showing how $0$-coin and $5$-coin deficits in $F(n,h,0)$ compensate for $2$-coin collisions in $T_1, T_2, T_5$.}
\label{tab:cancellation}
\end{table}

Notice in particular:
\begin{itemize}[nosep]
  \item For $k_0 = 5$, $F(n, 4, 0) \cup T_1(n, 5)$ and $F(n, 4, 0) \cup T_2(n, 5)$ are both subsets of $F(n, 6, 0)$ (\texttt{build\_T1\_subset\_llw\_F}, \texttt{build\_T2\_subset\_llw\_F}), which immediately bounds all diagonals $d \ne n - 1$ by $5$ via Lemma~\ref{lem:sharp-F}(d), leaving only the single diagonal $d = n - 1$ to be checked!
  \item For $k_0 = 7, n \equiv 3 \pmod 6$ ($n = 6k + 3$), $T_5(n)$ has \textbf{two distinct downward-diagonal collisions} when $k \ge 3$: one at $d = 6k + 6$ between row $2k$ and row $6k + 2 = n - 1$, and a second at $d = 6k - 1 = n - 4$ between row $4k - 2$ and row $6k - 1 = n - 4$. In \texttt{build\_T5\_dds\_ok\_two}, splitting $T_5(n)$ at the threshold $r < n - 4$ separates both pairs of colliding rows and proves $|T_5(n) \cap \{r + c = d\}| \le 2$ for all $d$, while on $d = n - 1 = 6k + 2$ (\texttt{build\_T5\_dds\_ok\_n\_sub\_one}) there is at most $1$ coin, so $|(F(n,6,0) \cup T_5(n)) \cap \{r + c = d\}| \le \max(6 + 1, 5 + 2) = 7$.
  \item Even on \emph{upward} diagonals, $T_5(n)$ has a $2$-coin collision in the corner region $d \ge 2n - 6$ (specifically at $d = 2n - 6$, which contains both $(n-4, 0)$ and $(n-1, 3)$, or $(n-2, 3)$ and $(n-1, 4)$ at $d = 2n - 6$). However, for $d \ge 2n - 6$, the entire $n \times n$ board has only $(2n - 1) - d \le 5 < 7$ cells on upward diagonal $d$ (\texttt{filter\_uds\_le\_board}), so the upward-diagonal bound $\le 7$ holds automatically there, while for $d < 2n - 6$, \texttt{build\_T5\_uds\_ok} proves $|T_5(n) \cap \{(n-1)+r-c = d\}| \le 1$.
\end{itemize}

\subsection{Gap 3: Explicit $c = 1$ ($n$-Queens) Proof}

For $c = 1$, \cite{LLW1990} simply cites the classical $n$-queens literature. To make the Lean~4 formalization completely self-contained, we formalized the explicit $n$-queens construction \texttt{classic\_n\_queens n} for all $n \ge 4$, splitting into the three standard modular cases:
\begin{itemize}[nosep]
  \item $n \not\equiv 2, 3 \pmod 6$ (\texttt{mem\_classic\_n\_queens\_other}): simple two-layer knight-step pattern $2r + 1$ for $0 \le r < n/2$ and $2(r - n/2)$ for $n/2 \le r < n$.
  \item $n = 6k + 2$ (\texttt{mem\_classic\_n\_queens\_mod6\_2}): perturbed pattern with special rows $r \in \{3k+1, 3k+2, 6k+1\}$ placed at columns $2, 0, 4$.
  \item $n = 6k + 3$ (\texttt{mem\_classic\_n\_queens\_mod6\_3}): perturbed pattern with special rows $r \in \{3k, 6k+1, 6k+2\}$ placed at columns $1, 0, 2$.
\end{itemize}
All five properties (\texttt{rows\_ok}, \texttt{cols\_ok}, \texttt{dds\_ok}, \texttt{uds\_ok}, \texttt{bc\_ok}) are proved by case analysis and Presburger arithmetic (\texttt{omega}).

\section{Conclusion / Referee Verdict on Le--Li--Wang (1990)}

From a mathematical and formal-verification standpoint:
\begin{enumerate}[nosep]
  \item \textbf{Is the LLW proof idea valid?} \textbf{Yes.} Their modular reduction $c = k_0 + 6t$ using slope-$2$ strips $F(n, h, s)$ and half-shifted transversals $T_1, T_2, T_4, T_5$ is genuinely sound and provides a polynomial-time explicit construction for every $(n, c)$.
  \item \textbf{Is the 1990 paper complete as written?} \textbf{No.} Treating~\cite{LLW1990} as a modern referee report, it is a proof sketch that omits the required complement construction at $c \in \{n - 2, n - 1\}$ (where their modular shift formula overflows and wraps around) and glosses over the $3 \mid n$ diagonal collisions in $T_1, T_2, T_5$ and their exact cancellation against the anti-diagonal deficits of $F(n, 2, 0)$ and $F(n, 6, 0)$.
  \item \textbf{Status in Lean~4:} With these completions in place, the theorem is now 100\% formally verified in Lean~4 (\texttt{n\_c\_coins\_placement\_proof.lean}).
\end{enumerate}

\begin{thebibliography}{9}
\bibitem{LLW1990}
Mao-Hua Le, Kang-Di Li, and Zhen-Fa Wang,
\emph{A note on the generalized $n$-queens problem},
Journal of Systems Science and Mathematical Sciences \textbf{10} (1990), no.~2, 167--168.
\url{https://sysmath.cjoe.ac.cn/jssc/CN/PDF/44564#page=2}

\bibitem{MO510135}
Luc Duponcheel,
\emph{Generalized $n$-Queens Conjecture (replacing the traditional $n$-queens problem with a coin placement problem)},
MathOverflow Question~510135 (2026).
\url{https://mathoverflow.net/questions/510135}
\end{thebibliography}

\end{document}
